% !TEX root = ../../main.tex
%% Défi du chapitre 6 — Fonctions plates et fonctions plateaux
%% Résultats vérifiés par verifications/defis_06_10.py.

\rubrique{Défi}

\begin{defi}{Fonctions plates et fonctions plateaux}{2 h}
On définit $f:\R\to\R$ par $f(x)=\mathrm{e}^{-1/x}$ si $x>0$ et $f(x)=0$ si
$x\leqslant0$. On va montrer que $f$ est indéfiniment dérivable, que
toutes ses dérivées sont nulles en $0$ — et pourtant $f$ n'est nulle sur
aucun voisinage à droite de $0$. On s'en servira pour raccorder des
fonctions de manière parfaitement lisse.

\partieExo{A}{Les dérivées sur $\intOO{0}{+\infty}$}
\begin{questions}
  \item Montrer que $f$ est continue sur $\R$.
  \item Montrer qu'il existe des polynômes $P_n$ tels que, pour tout
    $x>0$, $f^{(n)}(x)=P_n\bigl(\tfrac1x\bigr)\mathrm{e}^{-1/x}$, avec
    $P_0=1$ et $P_{n+1}=X^{2}\bigl(P_n-P_n'\bigr)$. Calculer $P_1$ et
    $P_2$.
  \item Montrer que $P_n$ est de degré $2n$, de coefficient dominant $1$.
\end{questions}

\partieExo{B}{Une fonction plate en $0$}
\begin{questions}[resume]
  \item Montrer que pour tout polynôme $P$,
    $P\bigl(\tfrac1x\bigr)\mathrm{e}^{-1/x}\xrightarrow[x\to0^{+}]{}0$.
  \item Montrer par récurrence que $f$ est de classe $\mathcal{C}^{n}$ sur
    $\R$ et que $f^{(n)}(0)=0$, pour tout $n\in\N$. \emph{On pourra utiliser
    le théorème de la limite de la dérivée.}
  \item En déduire que $f$ est de classe $\mathcal{C}^{\infty}$ sur $\R$,
    que $f(x)=o\bigl(x^{n}\bigr)$ en $0$ pour tout $n$, et pourtant que $f$
    n'est pas nulle au voisinage de $0$. Pourquoi ce phénomène est-il
    impossible pour une fonction polynomiale ?
\end{questions}

\partieExo{C}{Raccords lisses : fonctions plateaux}
\begin{questions}[resume]
  \item Soit $g(x)=f(x)\,f(1-x)$. Montrer que $g$ est de classe
    $\mathcal{C}^{\infty}$, strictement positive sur $\intOO{0}{1}$ et nulle
    ailleurs ; déterminer son maximum.
  \item Soit
    \[
    h(x)=\frac{f(x)}{f(x)+f(1-x)} .
    \]
    Justifier que $h$ est bien
    définie et de classe $\mathcal{C}^{\infty}$ sur $\R$, que $h=0$ sur
    $\intOF{-\infty}{0}$, $h=1$ sur $\intFO{1}{+\infty}$, et que $h$ est
    strictement croissante sur $\intFF{0}{1}$. Calculer
    $h(x)+h(1-x)$.
  \item Construire une fonction $\varphi$ de classe $\mathcal{C}^{\infty}$
    sur $\R$, égale à $1$ sur $\intFF{-1}{1}$, nulle en dehors de
    $\intOO{-2}{2}$, à valeurs dans $\intFF{0}{1}$ (une \emph{fonction
    plateau}).
\end{questions}

\medskip
\begin{center}
\begin{tikzpicture}[x=1.15cm,y=2.1cm,line cap=round]
  \draw[otRule,line width=0.3pt] (-2.8,1) -- (2.8,1);
  \draw[otMuted,line width=0.5pt,-{Stealth[length=1.6mm]}] (-2.9,0) -- (3,0);
  \draw[otMuted,line width=0.5pt,-{Stealth[length=1.6mm]}] (0,-0.12) -- (0,1.25);
  \foreach \v in {-2,-1,1,2}{\draw[otMuted,line width=0.4pt] (\v,0.03) -- (\v,-0.03);
    \node[font=\tiny,text=otMuted,below] at (\v,-0.03) {$\v$};}
  \draw[otGreenDark,line width=1.3pt] (-2.8,0) -- (-1.96,0)
    -- plot[domain=-1.96:-1.04,samples=120]
    ({\x},{exp(-1/(\x+2))/(exp(-1/(\x+2))+exp(-1/(-1-\x)))})
    -- (-1.04,1) -- (1.04,1)
    -- plot[domain=1.04:1.96,samples=120]
    ({\x},{exp(-1/(2-\x))/(exp(-1/(2-\x))+exp(-1/(\x-1)))})
    -- (1.96,0) -- (2.8,0);
  \draw[otNavy,line width=0.9pt,dash pattern=on 3pt off 1.6pt] (-2.8,0) -- (0.07,0)
    -- plot[domain=0.07:2.8,samples=120] ({\x},{exp(-1/\x)});
  \node[font=\tiny,text=otMuted,anchor=east] at (0,1) {$1$};
  \node[font=\scriptsize,text=otGreenDark,anchor=south] at (0,1.02) {$\varphi$};
  \node[font=\scriptsize,text=otNavy,anchor=west] at (2.3,0.62) {$f$};
\end{tikzpicture}\\[0.5mm]
\legende{En tirets, $f$ : nulle à gauche, elle quitte $0$ « sans
décoller » ; en vert, la fonction plateau $\varphi$ construite à la
question 9, de classe $\mathcal{C}^{\infty}$.}
\end{center}
\end{defi}
